\chapter{Introdução}
\label{chap:intro}

\section{Enquadramento e Motivação}
\label{sec:enquadramento}

Este relatório apresenta a segunda e última fase do projeto \textbf{GenJazz}, desenvolvido no âmbito da unidade curricular de \ac{IHC} do curso de Engenharia Informática da \ac{UBI}. O GenJazz é uma aplicação móvel que gera progressões harmónicas de jazz a partir de parâmetros configuráveis --- tonalidade, estrutura formal e tipo de modulação --- e permite ao utilizador ouvir o resultado, visualizá-lo numa grelha de compassos e guardá-lo numa biblioteca pessoal.

A ideia surgiu de um problema concreto: o jazz tem uma linguagem harmónica rica mas pouco acessível a quem está a aprender. Há livros e partituras, claro, mas faltam ferramentas interativas que permitam explorar essa linguagem de forma rápida e auditiva. O GenJazz tenta preencher esse espaço de forma simples e direta.

Na primeira fase foram definidos o conceito, as personas, os fluxos de navegação e o protótipo de baixa-média fidelidade. Nesta fase, o protótipo foi implementado como aplicação web funcional e melhorado com base nas heurísticas de Nielsen~\cite{nielsen1994} e em princípios de design inclusivo.

\section{Objetivos}
\label{sec:objetivos}

Os principais objetivos desta fase foram: implementar o protótipo da Parte~1 como aplicação web funcional; integrar a API GenJazz para geração de progressões e conversão para áudio; e adicionar funcionalidades de acessibilidade para utilizadores com deficiência de visão cromática que não tinham sido contempladas no protótipo original.

\section{Organização do Documento}
\label{sec:organizacao}

O documento está organizado da seguinte forma: o Capítulo~2 descreve a arquitetura do sistema; o Capítulo~3 é o manual do utilizador; o Capítulo~4 apresenta as decisões de design e acessibilidade; o Capítulo~5 detalha as alterações face ao protótipo da Parte~1; e o Capítulo~6 contém as conclusões e o que ficou por implementar. O Apêndice~A é a declaração de uso de IA.
